\chapter{Arthroplasty}

\section*{Introduction \& Clinical Indications}
Arthroplasty is a major orthopedic surgical procedure involving the reconstruction or replacement of a diseased or damaged joint to restore its anatomical integrity, alleviate debilitating pain, and significantly improve the patient's quality of life. This complex intervention is most frequently performed on large, weight-bearing joints such like the hip and knee, but its principles extend to the shoulder, ankle, elbow, and even smaller joints of the hand and foot. The clinical scenario typically involves end-stage degenerative joint disease, inflammatory arthropathy, or post-traumatic arthritis where the articular cartilage and underlying subchondral bone have been irreversibly compromised. The surgical objective is to excise the damaged joint surfaces and replace them with artificial components, known as prostheses, meticulously engineered to replicate the natural joint's biomechanics, stability, and range of motion. Its clinical significance lies in its profound ability to transform the lives of patients suffering from chronic joint pain and functional impairment, enabling a return to mobility and independence.

\begin{itemize}
  \item Joint Replacement Surgery
  \item Joint Reconstruction Surgery
  \item Prosthetic Joint Implantation
  \item Total Joint Arthroplasty (TJA)
  \item Total Hip Replacement (THR)
  \item Total Knee Replacement (TKR)
\end{itemize}

The fundamental principle of arthroplasty is to remove the damaged articular surfaces of a joint and replace them with smooth, durable prosthetic components that allow for pain-free motion. This involves precise surgical techniques to resect the diseased bone and cartilage, prepare the remaining bone for implant fixation, and then securely implant the artificial joint components. The prostheses typically consist of metal alloys (such as cobalt-chromium or titanium), high-density polyethylene, and sometimes ceramic materials, chosen for their biocompatibility, strength, and low friction properties.

The surgical rationale is to eliminate the source of pain, which is often bone-on-bone friction or inflammation from damaged cartilage, and to restore the biomechanical alignment and stability of the joint. Technical goals include achieving optimal implant positioning, balancing soft tissues (ligaments, tendons) around the joint to ensure stability and range of motion, and minimizing blood loss and surgical trauma. Fixation of the prosthetic components to the bone can be achieved through:
\begin{itemize}
  \item Cemented fixation: Using polymethyl methacrylate (PMMA) bone cement to bond the implant to the bone.
  \item Uncemented (press-fit) fixation: Relying on the ingrowth of bone into a porous-coated implant surface for biological fixation.
  \item Hybrid fixation: A combination of cemented and uncemented components within the same joint.
\end{itemize}
The choice of fixation method depends on patient factors such as age, bone quality, and surgeon preference.

Early attempts at joint reconstruction date back to the 19th century, utilizing materials like ivory, glass, and even interpositional soft tissues to create new joint surfaces. However, these efforts often yielded limited success due to material incompatibility and poor fixation. A significant milestone occurred in 1940 when Austin Moore performed the first metallic hip hemiarthroplasty, replacing only the femoral head with a cobalt-chromium alloy prosthesis. The modern era of total joint arthroplasty began in the 1960s with Sir John Charnley in the United Kingdom. Charnley pioneered the concept of low-friction arthroplasty for the hip, combining a small femoral head with a polyethylene acetabular cup, and crucially, introducing polymethyl methacrylate (PMMA) bone cement for secure fixation of both components. This innovation dramatically improved implant longevity and clinical outcomes, establishing the blueprint for contemporary total hip replacement. Subsequent decades saw advancements in materials science, surgical techniques, and implant designs, leading to the widespread success of total knee arthroplasty in the 1970s and the continuous evolution of arthroplasty for other joints.

\begin{itemize}
  \item Severe osteoarthritis, characterized by progressive cartilage loss, bone-on-bone friction, and debilitating pain, unresponsive to conservative management.
  \item Rheumatoid arthritis and other inflammatory arthropathies causing extensive joint destruction, pain, and deformity.
  \item Post-traumatic arthritis resulting from previous joint injuries, fractures, or dislocations that lead to premature degenerative changes.
  \item Avascular necrosis (osteonecrosis) of the femoral head or other joint surfaces, where bone tissue dies due to interrupted blood supply, leading to collapse and pain.
  \item Certain complex fractures, such as displaced femoral neck fractures in elderly patients, where arthroplasty offers better functional outcomes than internal fixation.
  \item Congenital joint deformities or developmental dysplasia that cause premature arthritis and significant functional limitations.
  \item Failed previous joint surgery, including failed osteotomies or prior arthroplasty with implant loosening, infection, or mechanical failure requiring revision.
  \item Benign or malignant bone tumors affecting the joint, necessitating resection and reconstruction.
\end{itemize}

Arthroplasty is predominantly considered a primary definitive treatment for end-stage joint arthritis and certain acute conditions like displaced femoral neck fractures in the elderly. It is typically indicated when non-surgical treatments, such as physical therapy, medications (NSAIDs, corticosteroids), activity modification, and injections, have failed to provide adequate pain relief or functional improvement. In this context, it serves as a primary intervention to restore joint function and significantly improve quality of life. However, arthroplasty can also serve as a secondary or salvage procedure. This is the case in revision arthroplasty, where a previously implanted prosthesis has failed due to loosening, infection, fracture, or wear, and the original components need to be removed and replaced. It may also be considered secondary in situations where less invasive surgical options, like arthroscopy or osteotomy, have been attempted but did not achieve long-term success.

\section*{Relevant Surgical Anatomy}
Successful arthroplasty necessitates an intimate understanding of the complex anatomy surrounding the target joint. For \textbf{total hip arthroplasty (THA)}, key structures include the acetabulum of the pelvis, the femoral head and neck, and the surrounding soft tissues. The acetabulum is a concave socket formed by the ilium, ischium, and pubis, articulating with the spherical femoral head. Important ligaments stabilizing the hip joint include the iliofemoral, pubofemoral, and ischiofemoral ligaments. Major muscle groups involved are the gluteal muscles (medius, minimus, maximus) for abduction and extension, adductors, and the iliopsoas for flexion. Critical neurovascular structures include the femoral nerve, artery, and vein located anteriorly, and the sciatic nerve positioned posteriorly, which is particularly vulnerable during posterior approaches. Surgical landmarks include the greater trochanter of the femur and the anterior superior iliac spine (ASIS). Lymphatic drainage from the hip region primarily flows to the superficial and deep inguinal lymph nodes.

For \textbf{total knee arthroplasty (TKA)}, the primary anatomical components are the distal femur, proximal tibia, and patella. The femoral condyles articulate with the tibial plateau, stabilized by the cruciate ligaments (anterior and posterior) and collateral ligaments (medial and lateral). The patella articulates with the trochlear groove of the femur. The quadriceps femoris muscle and patellar tendon are crucial for knee extension, while the hamstrings provide flexion. The popliteal fossa, located posteriorly, houses the vital popliteal artery and vein, and the tibial and common peroneal nerves, all of which are at risk during surgical dissection. Surgical landmarks include the patella, tibial tubercle, and epicondyles of the femur. Lymphatic drainage from the knee region primarily involves the popliteal lymph nodes.

\section*{Preoperative Workup \& Preparation}
Thorough pre-operative preparation is crucial to optimize patient outcomes and minimize risks. This typically begins several weeks before surgery with a comprehensive medical evaluation by the orthopedic surgeon and often by a primary care physician or an anesthesiologist. Key preparatory steps include:
\begin{itemize}
  \item \textbf{Medical Clearance:} Assessment of overall health, including cardiac, pulmonary, renal, and endocrine systems, to ensure the patient is fit for surgery and anesthesia. This may involve an electrocardiogram (ECG), chest X-ray, and pulmonary function tests.
  \item \textbf{Laboratory Tests:} Baseline blood work including complete blood count (CBC), electrolyte panel, kidney and liver function tests, coagulation profile (PT/INR, PTT), and blood type and cross-match in anticipation of potential transfusion.
  \item \textbf{Infection Screening:} Screening for and treatment of any active infections, particularly dental infections (dental clearance may be required), urinary tract infections, or skin infections, as these can significantly increase the risk of periprosthetic joint infection.
  \item \textbf{Medication Review:} Adjustment or temporary discontinuation of certain medications, especially anticoagulants (e.g., warfarin, novel oral anticoagulants) and antiplatelet agents (e.g., aspirin, clopidogrel), under medical supervision. Diabetic medications may also need adjustment.
  \item \textbf{Nutritional Optimization:} Encouragement of a healthy diet and, in some cases, nutritional supplements to support healing.
  \item \textbf{Pre-operative Education:} Patients receive detailed instructions regarding the surgical procedure, expected recovery, pain management, and post-operative precautions.
  \item \textbf{NPO Status:} Patients are instructed to be NPO (nil per os - nothing by mouth) for a specified period (typically 6-8 hours) before surgery to prevent aspiration during anesthesia.
  \item \textbf{Antibiotic Prophylaxis:} Administration of prophylactic antibiotics intravenously shortly before the incision to reduce the risk of surgical site infection.
  \item \textbf{Informed Consent:} Detailed discussion of the procedure, potential benefits, risks, and alternatives, culminating in the patient providing informed consent.
\end{itemize}

\textbf{Absolute Contraindications:}
\begin{itemize}
  \item Active local or systemic infection: This is the most critical contraindication, as it carries an extremely high risk of periprosthetic joint infection, which can be devastating.
  \item Severe neurotrophic arthropathy (e.g., Charcot joint): The lack of protective sensation and progressive joint destruction makes implant survival highly unlikely.
  \item Rapidly progressive neurological disease: Conditions that severely impair motor function or sensation, making rehabilitation impossible or leading to early implant failure.
  \item Inadequate soft tissue coverage: Insufficient healthy skin and muscle to cover the implant, increasing infection risk and wound complications.
  \item Skeletal immaturity: In growing children, arthroplasty can interfere with growth plates.
  \item Uncontrolled severe medical comorbidities: Such as unstable angina, recent myocardial infarction, severe uncontrolled diabetes, or severe pulmonary disease, which pose an unacceptably high anesthetic or surgical risk.
\end{itemize}

\textbf{Relative Contraindications and Precautions:}
\begin{itemize}
  \item Morbid obesity: Increases surgical complexity, anesthetic risks, wound complications, and implant loosening rates.
  \item Uncontrolled diabetes mellitus: Elevates infection risk and impairs wound healing. Strict glycemic control is essential.
  \item Severe osteoporosis: May compromise implant fixation and increase the risk of periprosthetic fracture.
  \item Severe peripheral vascular disease: Can impair wound healing and increase infection risk.
  \item Chronic osteomyelitis in the affected limb: Requires careful evaluation and often staged treatment to eradicate infection before arthroplasty.
  \item Young age: While not an absolute contraindication, younger, highly active patients may experience earlier implant wear and loosening, potentially requiring revision surgery.
  \item Psychiatric illness or cognitive impairment: May affect patient compliance with post-operative rehabilitation and precautions.
  \item History of deep vein thrombosis or pulmonary embolism: Requires careful thromboprophylaxis planning.
  \item Smoking: Increases risks of wound complications, infection, and impaired bone healing. Cessation is strongly advised pre-operatively.
\end{itemize}

\section*{Surgical Technique \& Procedure}
The surgical technique for arthroplasty is highly standardized yet tailored to the specific joint and patient needs. The procedure typically commences with the administration of anesthesia, which can range from general anesthesia to regional techniques such as spinal or epidural anesthesia, often supplemented with peripheral nerve blocks for enhanced postoperative pain control. The patient is then meticulously positioned on the operating table to optimize surgical exposure while maintaining physiological stability; for instance, hip arthroplasty may utilize a supine or lateral decubitus position, whereas knee arthroplasty is almost universally performed supine. A sterile surgical field is established, and a carefully planned incision is made through the skin and underlying soft tissues to gain access to the joint. The specific surgical approach (e.g., posterior, anterior, or anterolateral for the hip; medial parapatellar for the knee) is chosen based on surgeon preference, patient anatomy, and desired outcomes.

The key operative steps generally include:
\begin{enumerate}
  \item \textbf{Joint Exposure and Capsulotomy:} The joint capsule is incised and retracted to fully expose the damaged articular surfaces of the bone.
  \item \textbf{Bone Resection and Preparation:} Specialized surgical guides and oscillating saws are used to precisely resect the diseased cartilage and a thin layer of underlying subchondral bone, ensuring accurate alignment and component sizing. The remaining bone surfaces are then prepared by reaming, broaching, or drilling to create a precise fit for the prosthetic components.
  \item \textbf{Trial Reduction and Assessment:} Temporary trial components, made of plastic or metal, are inserted to assess joint stability, limb length, range of motion, and soft tissue tension. Adjustments are made to bone cuts or soft tissue releases as necessary to achieve optimal biomechanics.
  \item \textbf{Definitive Implant Placement:} The permanent prosthetic components, typically consisting of metal alloys and high-density polyethylene, are then implanted. These are secured to the bone using either polymethyl methacrylate (PMMA) bone cement, a press-fit technique relying on bone ingrowth, or a hybrid combination.
  \item \textbf{Soft Tissue Balancing and Closure:} Ligaments, tendons, and the joint capsule are meticulously balanced and repaired to ensure optimal joint stability and tracking throughout the full range of motion. The surgical wound is then closed in layers, including muscle, fascia, subcutaneous tissue, and skin, using sutures or staples.
  \item \textbf{Drain Placement (Optional):} A surgical drain may be temporarily placed within the joint space or subcutaneous tissue to evacuate any postoperative bleeding and reduce hematoma formation.
\end{enumerate}
A sterile dressing is applied, and the patient is transferred to the Post-Anesthesia Care Unit (PACU) for immediate postoperative monitoring.

\section*{Complications \& Risks}
While arthroplasty is generally safe and highly effective, like any major surgical procedure, it carries potential risks and complications. These can be broadly categorized as general surgical risks and procedure-specific risks.

\textbf{General Surgical and Anesthesia Risks:}
\begin{itemize}
  \item \textbf{Anesthesia Complications:} Reactions to anesthetic agents, respiratory depression, cardiac events (e.g., myocardial infarction, arrhythmia), stroke, nerve damage from positioning.
  \item \textbf{Bleeding:} Intraoperative blood loss requiring transfusion, post-operative hematoma formation.
  \item \textbf{Infection:} Superficial wound infection or, more seriously, deep periprosthetic joint infection (PJI), which can necessitate further surgery and prolonged antibiotic treatment.
  \item \textbf{Deep Vein Thrombosis (DVT) and Pulmonary Embolism (PE):} Blood clots forming in the leg veins that can travel to the lungs, a potentially life-threatening complication.
\end{itemize}

\textbf{Procedure-Specific Risks:}
\begin{itemize}
  \item \textbf{Joint Dislocation:} The prosthetic joint components can dislocate, particularly in the early post-operative period, often requiring reduction under anesthesia.
  \item \textbf{Limb Length Discrepancy:} The operated limb may be slightly longer or shorter than the unoperated limb, which can sometimes require shoe lifts or further intervention.
  \item \textbf{Nerve or Vascular Injury:} Damage to nerves (e.g., sciatic nerve in hip surgery, peroneal nerve in knee surgery) or blood vessels during the procedure, potentially leading to weakness, numbness, or circulatory problems.
  \item \textbf{Implant Loosening or Failure:} Over time, the prosthetic components can loosen from the bone or wear out, necessitating revision surgery.
  \item \textbf{Periprosthetic Fracture:} A fracture occurring around the implanted prosthesis, either during surgery or post-operatively due to trauma or stress.
  \item \textbf{Persistent Pain:} Despite successful surgery, some patients may experience ongoing pain, which can be due to various factors including nerve irritation, soft tissue impingement, or complex regional pain syndrome.
  \item \textbf{Joint Stiffness or Limited Range of Motion:} Scar tissue formation or inadequate rehabilitation can lead to a stiff joint, sometimes requiring manipulation under anesthesia.
  \item \textbf{Heterotopic Ossification:} Formation of new bone in soft tissues around the joint, which can restrict motion.
  \item \textbf{Allergic Reaction:} Rare reactions to the metal components of the prosthesis.
\end{itemize}

\section*{Postoperative Care \& Recovery}
Immediately after arthroplasty, the patient is transferred to the Post-Anesthesia Care Unit (PACU) for close monitoring of vital signs, pain levels, and surgical site. Pain management is a priority, often involving a multimodal approach with intravenous medications, nerve blocks, and oral analgesics. Early mobilization is encouraged, often within hours of surgery, to prevent complications like DVT and promote recovery. Physical therapists typically begin working with the patient on the first post-operative day, initiating exercises to restore range of motion and strength.

During the first 24-48 hours, patients are monitored for signs of bleeding, infection, or other acute complications. Drains, if placed, are usually removed within this period. Chemical thromboprophylaxis (e.g., anticoagulants) is continued to prevent blood clots. Patients are educated on weight-bearing restrictions, joint precautions (e.g., avoiding certain hip movements to prevent dislocation), and the use of assistive devices like walkers or crutches. Discharge from the hospital typically occurs within 1 to 3 days, depending on the patient's progress and the type of arthroplasty. Upon discharge, patients continue with prescribed pain medication, wound care instructions, and a detailed physical therapy regimen, either at home, in an outpatient clinic, or sometimes in a rehabilitation facility. The first 1-2 weeks focus on managing pain and swelling, protecting the surgical site, and gradually increasing mobility and strength. Long-term recovery involves continued physical therapy for several months, with a gradual return to normal activities and, eventually, low-impact sports.

During arthroplasty, the surgeon meticulously documents intraoperative findings, which typically include the extent of articular cartilage degeneration, presence of osteophytes (bone spurs), subchondral bone sclerosis or cysts, and the condition of the synovium. For osteoarthritis, findings often reveal full-thickness cartilage loss, eburnation (polished, ivory-like appearance) of the subchondral bone, and significant marginal osteophyte formation. In cases of rheumatoid arthritis, a hyperplastic, inflamed synovium with pannus formation eroding cartilage and bone is characteristic. Avascular necrosis presents with areas of collapsed, necrotic bone, often in the femoral head. Any resected tissue, such as portions of the femoral head, tibial plateau, menisci, or synovial tissue, is routinely sent for histopathological examination. The pathology report provides definitive confirmation of the underlying disease process, detailing the microscopic changes consistent with osteoarthritis, inflammatory arthritis, avascular necrosis, or, rarely, neoplastic processes. This pathological confirmation is vital for understanding the disease etiology and guiding any further medical management.

A normal or successful postoperative course following arthroplasty is characterized by several key indicators. Radiographically, the prosthetic components should appear anatomically aligned and well-fixed to the bone, without signs of loosening or malposition. Clinically, the patient should experience a progressive decrease in pain, leading to significant relief compared to their pre-operative state. The operated joint should demonstrate improved range of motion and stability on physical examination, allowing for functional activities. The surgical wound should heal cleanly, without signs of infection such as excessive redness, swelling, warmth, or purulent drainage. Furthermore, the patient should be able to participate effectively in physical therapy, gradually regaining strength and mobility, and progressing towards independent ambulation and a return to desired activities.

Post-operative care for arthroplasty involves a structured follow-up schedule to monitor recovery and detect potential complications.
\begin{itemize}
  \item \textbf{Orthopedic Clinic Visits:} Typically scheduled at 2 weeks (for wound check and suture/staple removal), 6 weeks, 3 months, 6 months, 1 year, and then annually or biennially for several years.
  \item \textbf{Imaging:} Follow-up X-rays of the operated joint are usually taken at key intervals (e.g., 6 weeks, 1 year, 5 years, 10 years) to monitor implant position, assess for signs of loosening, wear, or periprosthetic bone changes.
  \item \textbf{Physical Therapy (PT) / Occupational Therapy (OT):} Continued outpatient physical therapy is crucial for several weeks to months to maximize strength, flexibility, and functional recovery. Occupational therapy may be recommended to adapt daily activities.
  \item \textbf{Medication Management:} Review of pain medications, DVT prophylaxis, and any other necessary medications.
  \item \textbf{Activity Resumption Guidance:} Gradual progression of activities, with specific instructions on avoiding high-impact sports or activities that could jeopardize the implant.
  \item \textbf{Warning Signs Education:} Patients are educated on warning signs that necessitate immediate medical attention, such as fever, chills, increasing redness or drainage from the wound, sudden severe pain, or instability of the joint.
\end{itemize}
Regular follow-up is essential for the long-term success and surveillance of the prosthetic joint.

\section*{Outcomes \& Prognosis}
Arthroplasty procedures, particularly total hip arthroplasty (THA) and total knee arthroplasty (TKA), are among the most common and successful elective surgeries performed worldwide. The incidence of these procedures has been steadily increasing over the past few decades, driven by an aging population, rising rates of obesity, and expanding indications for surgery. In the United States, over 1 million total hip and knee arthroplasties are performed annually, with projections indicating continued growth. Osteoarthritis is by far the leading indication for both THA and TKA, accounting for approximately 90\% of cases. Rheumatoid arthritis, avascular necrosis, and post-traumatic arthritis represent smaller but significant proportions. While arthroplasty is performed across all adult age groups, the peak incidence is typically in individuals aged 60-80 years. There is a slight female predominance for TKA, while THA incidence is more evenly distributed between sexes. Mortality rates associated with elective primary arthroplasty are very low, generally less than 0.5\%, but morbidity, including complications like infection, DVT, and dislocation, remains a significant concern, affecting a small but notable percentage of patients.

The outcomes and prognosis following arthroplasty are overwhelmingly positive, making it one of the most successful interventions in modern medicine for improving quality of life. For both total hip and total knee arthroplasty, success rates in terms of pain relief and functional improvement exceed 90-95\%. Patients typically experience a dramatic reduction in pain, improved mobility, and the ability to return to many daily activities that were previously limited. Implant survival rates are excellent, with 10-year survival often exceeding 90-95\% and 20-year survival rates ranging from 70-85\% for primary THA and TKA. Prognostic factors influencing long-term success include patient age (younger, more active patients may have earlier wear), bone quality, adherence to rehabilitation protocols, and the absence of complications like infection or loosening. While recurrence of the original condition is prevented by joint replacement, complications such as periprosthetic joint infection, aseptic loosening, or dislocation can necessitate revision surgery. However, even in these cases, subsequent revision procedures often yield good results, albeit with potentially higher risks and less predictable outcomes than primary surgery.

\section*{References}
\begin{enumerate}
  \item MedlinePlus. Arthroplasty. U.S. National Library of Medicine; 2024. Available from: \url{https://medlineplus.gov/arthroplasty.html}
  \item Azar FM, Beaty JH. Campbell's Operative Orthopaedics. 14th ed. Philadelphia, PA: Elsevier; 2021.
  \item American Academy of Orthopaedic Surgeons (AAOS). Total Joint Replacement. Available from: \url{https://orthoinfo.aaos.org/en/treatment/total-joint-replacement/}
  \item Townsend CM Jr, Beauchamp RD, Evers BM, Mattox KL. Sabiston Textbook of Surgery: The Biological Basis of Modern Surgical Practice. 21st ed. Philadelphia, PA: Elsevier; 2021.
\end{enumerate}

\clearpage
